\documentclass[12pt,a4paper]{article}

\usepackage[margin=2.5cm]{geometry}
\usepackage{amsmath,amssymb,amsthm,amsfonts}
\usepackage{mathtools}
\usepackage{bm}
\usepackage{xcolor}
\usepackage[colorlinks=true,linkcolor=black,citecolor=black,urlcolor=blue]{hyperref}
\usepackage{booktabs}
\usepackage{array}
\usepackage{longtable}
\usepackage{enumitem}
\usepackage{microtype}
\usepackage{cleveref}
\usepackage{mathrsfs}
\usepackage{thmtools}
\usepackage{thm-restate}
\usepackage{graphicx}
\usepackage{multirow}

\theoremstyle{plain}
\newtheorem{theorem}{Theorem}[section]
\newtheorem{proposition}[theorem]{Proposition}
\newtheorem{lemma}[theorem]{Lemma}
\newtheorem{corollary}[theorem]{Corollary}

\theoremstyle{definition}
\newtheorem{definition}[theorem]{Definition}
\newtheorem{construction}[theorem]{Construction}

\theoremstyle{remark}
\newtheorem{remark}[theorem]{Remark}

\newcommand{\GFtwo}{\mathbb{F}_2}
\newcommand{\Z}{\mathbb{Z}}
\newcommand{\R}{\mathbb{R}}
\newcommand{\HX}{H_X}
\newcommand{\HZ}{H_Z}
\newcommand{\rank}{\operatorname{rank}}
\newcommand{\wt}{\operatorname{wt}}
\newcommand{\rowsp}{\operatorname{rowsp}}
\newcommand{\CD}{\mathscr{D}}
\newcommand{\Lat}{\mathscr{L}}
\newcommand{\FS}{\mathcal{F}}
\newcommand{\PS}{\mathcal{P}}

\definecolor{nogocolor}{rgb}{0.55,0.0,0.0}
\newcommand{\NOGO}{\textcolor{nogocolor}{\textbf{[No-Go]}}}
\definecolor{newcolor}{rgb}{0.0,0.35,0.15}
\newcommand{\NEW}{\textcolor{newcolor}{\textbf{[New]}}}

\title{\Large\bfseries A Lorentzian CSS Duality in Causal Diamond\\[4pt]
Quantum Error-Correcting Codes:\\[4pt]
Four Codes from One Geometry via\\[4pt]
Orientation Reversal}

\author{Yannick Schmitt \\ 
\texttt{ynnk.schmitt@gmail.com} \\
\vspace{4pt} 
\href{https://doi.org/10.5281/zenodo.19343889}{DOI: 10.5281/zenodo.19343889}}
\date{\today}

\begin{document}
\maketitle

\begin{abstract}
We show that the discrete Lorentzian causal diamond $\CD$ — a two-complex
built on the twelve lightlike nearest-neighbour vectors of the ternary
Minkowski lattice $\Lat = \{-1,0,+1\}^4$ — generates not two but
\emph{four} CSS quantum error-correcting codes via a geometric duality,
and that this duality fully explains the asymmetry $d_Z = 2$ that limited
the original two-code family.

\vspace{1em}\noindent\textbf{The duality map.}
Exchanging the roles of the 21 null plaquettes and the three
spatial-axis groups under the CSS transpose
$(H_X,\, H_Z) \mapsto (H_Z^{(0)},\, H_X^{(0)})$ — where $H_Z^{(0)} = M^T$
denotes the plaquettes used as X-type checks — produces a
\emph{dual family} that corrects Z-errors with distance 6 while
detecting X-errors.
Combining one primal and one dual measurement round on the same 21
plaquettes (in Z- and X-basis respectively) yields a two-stage protocol
that corrects both error types simultaneously.

\vspace{1em}\noindent\textbf{Four codes.}
\emph{Code~I} $[[12,4,(4,2)]]$ encodes 4 logical qubits at rate $\tfrac{1}{3}$
and corrects all X-errors but only detects Z-errors ($d_Z=2$).
\emph{Code~II} $[[12,1,(4,3)]]$ is balanced and distance-3.
\emph{Dual~A} $[[12,2,(2,6)]]$ (new) corrects Z-errors with distance 6
but only detects X-errors ($d_X=2$); it uses the same 21 plaquettes as
X-type stabilisers with the three spatial-axis groups as Z-type stabilisers.
\emph{Dual~II} $[[12,1,(3,4)]]$ (new) exchanges the role of distances
relative to Code~II: $d_Z=4$ at the cost of $d_X=3$.

\vspace{1em}\noindent\textbf{Geometric origin of the asymmetry.}
The $d_Z=2$ obstruction in Code~I is proven to arise from a rank gap:
$\rank(M)=8$ over $\R$ but $\rank(M)=7$ over $\GFtwo$.
The extra GF(2)-flat direction is the all-ones vector,
which is a Z-logical in Code~I but a Z-stabiliser in Dual~A.
This gap is a consequence of the Lorentzian temporal boundary charge
$n_{\mathrm{eff}}^0 = 12$ computed in~\cite{SchmittCD}; the duality
corresponds precisely to the Euclidean (Riemannian) orientation in which
this charge vanishes.

\vspace{1em}\noindent\textbf{Numerical results.}
At $p=0.01$ under code-capacity depolarising noise (20000 trials),
Code~I achieves $P_{\log}(Z)=0.101$ (Z-detection fails),
Code~II achieves $P_{\log}=0.005$, and the two-stage combined protocol
achieves $P_{\log}=0.006$ with $k_{\mathrm{eff}}=2$ logical qubits.
Dual~II achieves $P_{\log}(Z)=0.001$ at $p=0.01$ (all weight-1
Z-errors corrected, $d_Z=4$).
The circuit-level threshold for Code~II is $p_c \approx 3.5\%$.
\end{abstract}

\tableofcontents
\bigskip

\noindent\textbf{MSC 2020:} 81P70, 94B05, 52B11.\\[2pt]
\textbf{Keywords:} causal diamond; Lorentzian CSS duality; orientation
reversal; asymmetric quantum code; distance asymmetry; D4 root system;
rank gap; two-stage protocol.

%------------------------------------------------------------------
\section{Introduction}\label{sec:intro}
%------------------------------------------------------------------

\subsection{Overview and motivation}

The connection between discrete spacetime geometry and quantum
error-correcting codes is well established: the surface code and toric
code arise from the homology of a torus~\cite{Fowler2012}, and hyperbolic
codes exploit negative curvature to achieve better rate--distance
tradeoffs~\cite{Panteleev2022,Leverrier2022}.
The present work shows that the discrete Lorentzian causal diamond
$\CD$, studied in the companion paper~\cite{SchmittCD}, carries a
geometric duality that generates a complete family of four CSS codes
and makes the distance asymmetry of the original construction
algebraically inevitable rather than merely observed.

The key insight is that the causal diamond supports two natural
orientations of its boundary: the \emph{Lorentzian} orientation, in
which past links are incoming, produces the boundary charge
$n_{\mathrm{eff}}^0 = 12$ that motivates the all-ones X-check and
yields Code~I; the \emph{Euclidean} orientation, in which all links
are outgoing and the boundary sum vanishes, has no global check and
yields a dual family in which the 21 plaquettes serve as X-type checks
and the spatial-axis groups serve as Z-type checks.
The two orientations are related by the substitution $t\to ix$ (Wick
rotation), so the duality between primal and dual code families is a
QEC realisation of the Wick rotation.

\subsection{Summary of results}

\paragraph{Primal family} ($H_Z = M^T$, the 21-plaquette matrix).
\begin{itemize}[noitemsep]
\item \textbf{Code~I} $[[12,4,(4,2)]]$: all-ones X-check, rate $\tfrac{1}{3}$,
  parallel syndrome extraction, corrects X-errors, \emph{detects but
  cannot correct} Z-errors.
\item \textbf{Code~II} $[[12,1,(4,3)]]$: four weight-6 X-stabilisers
  assigning each qubit a unique 4-bit label, balanced distances.
\end{itemize}

\paragraph{Dual family} \NEW\ ($H_X = M^T$, the 21-plaquette matrix
used as X-type checks).
\begin{itemize}[noitemsep]
\item \textbf{Dual~A} $[[12,2,(2,6)]]$: three spatial-axis groups as
  Z-stabilisers, $d_Z = 6$, corrects all weight-1 and weight-2
  Z-errors, detects but cannot correct X-errors.
\item \textbf{Dual~II} $[[12,1,(3,4)]]$: Code~II's X-rows promoted to
  Z-checks, $d_Z = 4$, $d_X = 3$; strictly better than Code~II for
  dephasing-dominated hardware.
\end{itemize}

\paragraph{Two-stage combined protocol} \NEW.
Measuring the 21 plaquettes first in Z-basis (Stage~1, corrects
X-errors, $d_X=4$) and then in X-basis (Stage~2, corrects Z-errors,
$d_Z=6$) yields a protocol on $k_{\mathrm{eff}}=2$ logical qubits
with $P_{\log}=0.006$ at $p=0.01$.

\paragraph{No-go theorems (unchanged from v3).}
The Pigeonhole No-Go Theorem (\cref{thm:nogo}) proves that
$d_Z\geq3$ with $k\geq2$ is impossible in any CSS code on this
geometry, and Code~II saturates the bound.
The X-Decoration Equivalence (\cref{thm:nodeco}) extends this to
weight-$\leq6$ non-CSS codes.

\subsection{Comparison with existing codes}

\begin{table}[h]
\centering
\caption{All four causal diamond CSS codes alongside nearest-scale
  comparison codes. New codes marked \NEW.}
\label{tab:comparison}
\small
\begin{tabular}{lcccccl}
\toprule
Code & $n$ & $k$ & $(d_X,d_Z)$ & Rate & Max wt & Note \\
\midrule
Surface $[[9,1,3]]$   & 9  & 1  & $(3,3)$ & $\tfrac{1}{9}$  & 4 & planar \\
Steane $[[7,1,3]]$    & 7  & 1  & $(3,3)$ & $\tfrac{1}{7}$  & 4 & transversal \\
Toric $[[18,2,3]]$    & 18 & 2  & $(3,3)$ & $\tfrac{1}{9}$  & 4 & non-planar \\
\midrule
\textbf{Code~I}   $[[12,4,(4,2)]]$  & 12 & 4 & $(4,2)$ & $\mathbf{\tfrac{1}{3}}$ & 4$^*$ & primal, rate record \\
\textbf{Code~II}  $[[12,1,(4,3)]]$  & 12 & 1 & $(4,3)$ & $\tfrac{1}{12}$ & 6 & primal, balanced \\
\textbf{Dual~A}   $[[12,2,(2,6)]]$  & 12 & 2 & $(2,6)$ & $\tfrac{1}{6}$ & 4 & \NEW, dual, $d_Z=6$ \\
\textbf{Dual~II}  $[[12,1,(3,4)]]$  & 12 & 1 & $(3,4)$ & $\tfrac{1}{12}$ & 6 & \NEW, dual, $d_Z=4$ \\
\bottomrule
\end{tabular}
\end{table}

\noindent${}^*$After decomposing into three disjoint weight-4 checks.

%------------------------------------------------------------------
\section{The Causal Diamond Two-Complex}\label{sec:geometry}
%------------------------------------------------------------------

We recall the relevant facts from~\cite{SchmittCD}.

The \emph{ternary Minkowski lattice} is
\[
\Lat = \{-1,0,+1\}^4, \quad \eta = \mathrm{diag}(-1,+1,+1,+1).
\]
It contains exactly 12 lightlike non-origin vectors
\[
\mathcal{L} = \{(\pm1,\, \sigma_j\bm{e}_j) : j\in\{1,2,3\},\,\sigma_j\in\{\pm1\}\},
\]
partitioned as $\mathcal{L} = \FS \cup \PS$ with $|\FS|=|\PS|=6$.

\begin{definition}[Causal Diamond Two-Complex]
$\CD$ has 0-cells: origin and $\mathcal{L}$ (13 vertices);
1-cells: 12 null links $\{[\bm{0},L] : L\in\mathcal{L}\}$;
2-cells: 21 order-four plaquettes of type $2\FS+2\PS$.
\end{definition}

The link--plaquette incidence matrix $M\in\{0,1\}^{12\times21}$ has
$M_{lp}=1$ iff link $l$ belongs to plaquette $p$.
The plaquette Laplacian $K = MM^T$ has exact integer spectrum
\[
\operatorname{spec}(K) = \{0^{(4)},\; 6^{(2)},\; 8^{(3)},\; 10^{(2)},\; 28^{(1)}\};
\]
each link participates in exactly 7 plaquettes ($K_{ll}=7$).

\paragraph{Rank over $\R$ and $\GFtwo$.}
$\rank(M) = 8$ over $\R$ and $\rank(M) = 7$ over $\GFtwo$.
The rank gap of 1 is the central algebraic fact of this paper; its QEC
meaning is established in \cref{sec:rankgap}.

\paragraph{Spatial-axis decomposition.}
We partition the 12 null links into three groups by spatial axis:
$G_j = \{l \in \mathcal{L} : v_l^j \neq 0\}$ for $j = 1,2,3$.
Explicitly: $G_1 = \{0,5,6,11\}$, $G_2 = \{1,4,7,10\}$, $G_3 = \{2,3,8,9\}$.
The three sets are pairwise disjoint and partition $\{0,\ldots,11\}$.

%------------------------------------------------------------------
\section{Shared Code Structure and the GF(2) Rank Gap}\label{sec:shared}
%------------------------------------------------------------------

\begin{construction}[Z-stabiliser matrix]\label{con:HZ}
$\HZ := M^T \in \GFtwo^{21\times12}$.
Each row corresponds to one plaquette and has weight 4.
\end{construction}

\begin{proposition}\label{prop:HZ_rank}
$\rank(\HZ) = 7$ over $\GFtwo$.
The kernel $\ker(\HZ)$ has dimension 5 and contains exactly 27
non-zero vectors of weight 4 or 6 (3 of weight 4, 24 of weight 6).
\end{proposition}

\begin{proposition}\label{prop:syndrome}
All 12 single-qubit Z-error syndromes under $\HZ$ are distinct,
each of weight 7, with minimum Hamming distance 4 between any two.
\end{proposition}
\begin{proof}
Each qubit $l$ has syndrome weight 7 (from $K_{ll}=7$).
For $l \neq l'$: $d_H \geq 7+7-2|\text{shared plaquettes}| \geq 4$
since any two links share at most 5 plaquettes.
\end{proof}

\subsection{The rank-gap mechanism}\label{sec:rankgap}

The following theorem gives a geometric derivation of the constraint
$d_Z = 2$ in Code~I that was previously only observed numerically.

\begin{theorem}[GF(2) Rank Gap and the Temporal Obstruction]
\label{thm:rankgap}
\NEW\; The rank gap $\rank(M)_\R - \rank(M)_{\GFtwo} = 1$ has the
following algebraic cause and QEC consequence.

\emph{(i) Algebraic cause.}
The all-ones vector $\bm{1}_{12}$ lies in $\ker(\HZ)$ over $\GFtwo$
(every plaquette has weight 4, which is even) but does \emph{not} lie
in $\ker(\HZ)$ over $\R$ (the holonomy of every plaquette is 4, not 0).
This is the unique extra GF(2)-flat direction; it arises from the
ternary structure of $\Lat$.

\emph{(ii) Geometric cause.}
By Theorem~4.2 of~\cite{SchmittCD}, the Lorentzian boundary sum satisfies
$n_{\mathrm{eff}}^0 = 2|\FS| = 12$.
The weight of the all-ones vector equals this temporal charge.
The Euclidean orientation (all links outgoing) gives
$n_{\mathrm{eff},\mathrm{Eucl}}^0 = \sum_{l\in\FS} v_l^0 + \sum_{l\in\PS} v_l^0 = 6-6 = 0$,
eliminating the charge and with it the all-ones flat direction.

\emph{(iii) QEC consequence.}
In any CSS code with $H_X$ having only even-weight rows,
$\bm{1}_{12} \in \ker(H_X)$.
Since $\bm{1}_{12}$ is not in $\rowsp(\HZ)$ (minimum stabiliser weight
is 4), and any weight-2 vector $e_{q_1}+e_{q_2}$ satisfies the same
conditions, all 66 qubit pairs are Z-logicals.
Therefore $d_Z = 2$ and Code~I cannot correct Z-errors.
\end{theorem}
\begin{proof}
Part~(i): $\HZ \bm{1}_{12} = (4,4,\ldots,4)^T \equiv \bm{0}$ over $\GFtwo$
since every plaquette has exactly 4 members.
Over $\R$, the same product gives $(4,4,\ldots,4)^T \neq \bm{0}$.
The uniqueness of this direction follows from $\rank(M)_\R - \rank(M)_{\GFtwo} = 1$.
Parts~(ii) and (iii) follow directly from the definitions.
\end{proof}

\begin{remark}
\cref{thm:rankgap} makes precise why Code~I ``catches all Z-errors
but corrects none'': the single all-ones X-check detects any Z-error
(because every qubit is in the support of $H_X^{(I)}$) but provides
only 1 bit of syndrome information, which is insufficient to locate
the error. The dual codes (\cref{sec:dual}) resolve this by
contributing 7 independent bits of Z-location information.
\end{remark}

%------------------------------------------------------------------
\section{The Lorentzian CSS Duality}\label{sec:duality}
%------------------------------------------------------------------

\subsection{Statement of the duality}

\begin{theorem}[Lorentzian CSS Duality]\label{thm:duality}
\NEW\; Define the following two families of CSS codes on the same
12 physical qubits:

\begin{description}[style=nextline,labelindent=0pt]
\item[Primal family] $H_Z^{(0)} := M^T$ (21 plaquettes as Z-type checks),
  with $H_X$ chosen from among the valid X-check matrices.
\item[Dual family] $H_X^{(0)} := M^T$ (21 plaquettes as X-type checks),
  with $H_Z$ chosen from among the valid Z-check matrices.
\end{description}

The map
\[
\Phi:\; (H_X,\, H_Z^{(0)}) \;\longmapsto\; (H_X^{(0)},\, H_X^T)
\]
is a \emph{CSS duality}: if $(H_X, H_Z^{(0)})$ is a valid CSS code,
so is $\Phi(H_X, H_Z^{(0)})$.
Under $\Phi$:
\begin{itemize}[noitemsep]
\item X-logicals in the primal code become Z-logicals in the dual code.
\item Z-logicals in the primal code become X-logicals in the dual code.
\item $(d_X, d_Z)$ in the primal maps to $(d_Z^{(0)}, d_X^{(0)})$
  in the dual, where the superscript $(0)$ denotes the dual distances.
\end{itemize}

Applying $\Phi$ to Code~I and Code~II yields:
\begin{align*}
\Phi(\text{Code~I}) &= \text{Dual~A}: && [[12,\,2,\,(2,6)]], \\
\Phi(\text{Code~II}) &= \text{Dual~II}: && [[12,\,1,\,(3,4)]].
\end{align*}
\end{theorem}

\subsection{Geometric interpretation}

The duality $\Phi$ corresponds to changing the orientation of the
causal diamond boundary:

\begin{itemize}[noitemsep]
\item \textbf{Lorentzian orientation} (primal): future links outgoing,
  past links incoming. Temporal boundary sum $n_{\mathrm{eff}}^0 = 12$,
  spatial sums $n_{\mathrm{eff}}^j = 0$.
  The temporal charge motivates the all-ones Z-detection check;
  spatial cancellation gives the disjoint $G_j$ groups.
  \emph{Plaquettes as Z-type checks detect X-errors.}
\item \textbf{Euclidean orientation} (dual): all links outgoing.
  Total boundary sum $n_{\mathrm{eff},\mathrm{Eucl}} = (0,0,0,0)$.
  No global check survives; the spatial structure ($G_j$ as
  Z-stabilisers) dominates.
  \emph{Plaquettes as X-type checks detect Z-errors.}
\end{itemize}

This is a QEC realisation of the Wick rotation $t \to ix$: the
substitution that sends Minkowski signature to Euclidean signature
exchanges the two code families.

\subsection{Isometric syndrome structure}

A striking consequence of the duality is that the syndrome information
for X-errors and Z-errors is \emph{geometrically identical}:

\begin{proposition}[Isometric Syndrome Structure]\label{prop:isometric}
\NEW\; Every qubit $l$ participates in exactly 7 of the 21 plaquettes
(since $K_{ll}=7$). Therefore:
\begin{itemize}[noitemsep]
\item In the primal family: each qubit triggers 7 of the 21 Z-checks,
  giving a weight-7 X-error syndrome. All 12 single-qubit X-error
  syndromes are distinct.
\item In the dual family: each qubit triggers 7 of the 21 X-checks,
  giving a weight-7 Z-error syndrome. All 12 single-qubit Z-error
  syndromes are distinct.
\end{itemize}
In both cases, 7 independent bits of location information are available.
The geometry provides equally rich syndrome information for both error
types; the two code families each exploit only one half of this capacity.
\end{proposition}

%------------------------------------------------------------------
\section{Code~I: The Asymmetric High-Rate Code}\label{sec:codeI}
%------------------------------------------------------------------

\subsection{Construction and parameters}

\begin{construction}[Code~I]\label{con:codeI}
$H_X^{(I)} := \bm{1}_{1\times12}$ (the all-ones row vector).
\end{construction}

\begin{proposition}
$H_X^{(I)} H_Z^T = \bm{0}$ over $\GFtwo$: each plaquette has weight 4.
\end{proposition}

\begin{theorem}[Parameters of Code~I]\label{thm:codeI}
Code~I has parameters $[[12, 4, (d_X,d_Z)=(4,2)]]$ with rate $\tfrac{1}{3}$.
\end{theorem}
\begin{proof}
$k = 12 - 1 - 7 = 4$.
$d_X = 4$: the three group characteristic vectors $\bm{1}_{G_j}$ are
weight-4 X-logicals; no weight-1/2/3 X-logical exists (exhaustive check).
$d_Z = 2$: every weight-2 even vector lies in $\ker(H_X^{(I)})$ and
none lies in $\rowsp(\HZ)$ (minimum stabiliser weight is 4).
By \cref{thm:rankgap}, this is forced by the GF(2) rank gap.
\end{proof}

\begin{remark}[Why $d_Z = 2$ is inevitable for Code~I]
\cref{thm:rankgap}(iii) gives a geometric proof that no choice of
single-row even-weight $H_X$ on this geometry can achieve $d_Z > 2$.
The 31 other valid single-row CSS X-matrices all share this limitation.
The only remedy is to add more X-rows (Code~II) or switch to the dual
orientation (Dual~A).
\end{remark}

\subsection{Disjoint weight-4 X-stabiliser decomposition}\label{sec:xstab_A}

\begin{theorem}[Disjoint spanning triple]\label{thm:triple}
The three group characteristic vectors $r_j := \bm{1}_{G_j}$ satisfy:
(i)~$\wt(r_j)=4$;\;
(ii)~$\HZ r_j = 0$;\;
(iii)~supports pairwise disjoint;\;
(iv)~$r_1+r_2+r_3 = \bm{1}_{12}$ mod~2.
This is the unique such decomposition.
\end{theorem}
\begin{proof}
(ii): by the $2\FS+2\PS$ plaquette structure, each $G_j$ contributes
exactly 2 links to every plaquette, giving even intersection.
Exhaustive search over all $\binom{12}{4}=495$ weight-4 vectors finds
exactly 3 members of $\ker(\HZ)$; the triple $\{r_1,r_2,r_3\}$ is
thus uniquely determined.
(iv): $G_1,G_2,G_3$ partition $\{0,\ldots,11\}$.
\end{proof}

\begin{corollary}[Parallel syndrome extraction]\label{cor:parallel}
The X-syndrome can be extracted by three simultaneous weight-4 ancilla
circuits (one per group), reducing circuit depth to~1.
\end{corollary}

\subsection{X-error correction}

\begin{proposition}\label{prop:bp}
The 21-row Z-stabiliser matrix $\HZ$ admits a lookup decoder correcting
all weight-1 X-errors: the 12 single-qubit X-error syndromes are distinct
(\cref{prop:syndrome}), enabling unique identification.
At $p=0.05$ (independent bit-flip, 1000 trials) the logical X-error
rate is $\approx 4.6\%$.
\end{proposition}

\subsection{Code-capacity Z-error performance}

\begin{proposition}\label{prop:codeI_z}
Under independent depolarising noise (20000 trials), Code~I achieves:
\begin{center}
\small
\begin{tabular}{ccc}
\toprule
$p$ & $P_{\log}(X)$ & $P_{\log}(Z)$ \\
\midrule
0.001 & 0.0000 & 0.0102 \\
0.005 & 0.0003 & 0.0545 \\
0.010 & 0.0040 & 0.1011 \\
0.050 & 0.0383 & 0.4311 \\
\bottomrule
\end{tabular}
\end{center}
The Z-threshold is $p_c(Z) \approx 6.1\%$; no X-threshold is reached
below $p=0.20$.
As expected from $d_Z = 2$, the Z-logical error rate is substantial
at hardware-relevant error rates.
Code~I is therefore a Z-error-\emph{detecting} code; it requires
post-selection or erasure conversion for autonomous Z-correction,
or pairing with Dual~A in the two-stage protocol of \cref{sec:twostage}.
\end{proposition}

%------------------------------------------------------------------
\section{Code~II: The Balanced Distance-3 Code}\label{sec:codeII}
%------------------------------------------------------------------

\subsection{Construction principle}

For $d_Z \geq 3$ we need $H_X$ to separate every qubit pair:
each pair $(q_1,q_2)$ must have at least one X-check containing
exactly one of the two.
Together with single-qubit detection, this requires:
(i)~all 12 qubit labels distinct;
(ii)~no qubit has label $\bm{0}$;
(iii)~each X-check row in $\ker(\HZ)$.

\begin{theorem}[Code~II exists]\label{thm:codeII_exists}
There are exactly 1248 sets of 4 vectors in $\ker(\HZ)$ satisfying
(i)--(iii), all consisting of four weight-6 rows.
\end{theorem}
\begin{proof}
Exhaustive search over all $\binom{27}{4}$ 4-element subsets of the
27 non-zero weight-4/6 vectors in $\ker(\HZ)$.
The 3 weight-4 members cannot form a 4-row set (verified exhaustively).
\end{proof}

\begin{construction}[Code~II]\label{con:codeII}
Let $H_X^{(II)}$ be any of the 1248 valid 4-row matrices.
\end{construction}

\begin{theorem}[Parameters of Code~II]\label{thm:codeII}
Code~II has parameters $[[12, 1, (d_X,d_Z)=(4,3)]]$.
\end{theorem}
\begin{proof}
$k = 12 - 4 - 7 = 1$.
$d_Z = 3$: conditions~(i)--(ii) eliminate all weight-1 and weight-2
Z-logicals; exhaustive search confirms 16 weight-3 Z-logicals.
$d_X = 4$: the three group vectors $\bm{1}_{G_j}$ are X-logicals.
\end{proof}

\subsection{Role reversal of the all-ones vector}

\begin{proposition}\label{prop:all_ones_role}
The all-ones vector $\bm{1}_{12}$ is an X-\emph{stabiliser} in Code~I
and a weight-12 X-\emph{logical operator} in Code~II.
In the dual family, $\bm{1}_{12}$ becomes a Z-logical in Dual~A (it
lies in $\ker(H_X^{(0)} = M^T)$ by parity, and is not in
$\rowsp(H_Z^{(A)} = H_{\text{spatial}})$ since the minimum Z-stabiliser
weight in Dual~A is 4).
\end{proposition}

%------------------------------------------------------------------
\section{Dual Code Family}\label{sec:dual}
%------------------------------------------------------------------

\subsection{Dual~A: The Z-Correcting Dual}

\begin{construction}[Dual~A]\label{con:dualA}
\NEW\;
$H_X^{(A)} := M^T \in \GFtwo^{21\times12}$ (21 plaquettes as X-type checks).
$H_Z^{(A)} := H_{\text{spatial}} \in \GFtwo^{3\times12}$, where row~$j$
is the indicator of $G_j$.
\end{construction}

\begin{theorem}[CSS commutativity of Dual~A]\label{thm:dualA_comm}
$H_X^{(A)} (H_Z^{(A)})^T = \bm{0}$ over $\GFtwo$.
\end{theorem}
\begin{proof}
Entry $(p,j)$ of the product is $|p \cap G_j| \bmod 2$.
Every plaquette has type $2\FS+2\PS$; each future (past) link has
exactly one nonzero spatial component. Therefore every plaquette
contributes exactly 0 or 2 links to each spatial axis group $G_j$,
giving $|p \cap G_j| \in \{0,2\}$, which is always even. 
\end{proof}

\begin{theorem}[Parameters of Dual~A]\label{thm:dualA}
\NEW\; Dual~A has parameters $[[12, 2, (d_X,d_Z)=(2,6)]]$.
\end{theorem}
\begin{proof}
$k = 12 - \rank(H_X^{(A)}) - \rank(H_Z^{(A)}) = 12 - 7 - 3 = 2$.

$d_Z = 6$:
The G-j groups, which were weight-4 X-logicals in Code~I, are now
Z-\emph{stabilisers} in $\rowsp(H_Z^{(A)})$.
Their promotion removes all weight-4 Z-logicals from the logical space.
The 24 weight-6 vectors in $\ker(M^T)$ that lie in $\ker(H_X^{(A)})$
and are not in $\rowsp(H_Z^{(A)})$ form the minimum-weight Z-logical set.
Exhaustive check confirms 24 weight-6 Z-logicals and no lighter ones.

$d_X = 2$:
Any weight-2 vector $e_{q_1}+e_{q_2}$ with $q_1,q_2$ in the same $G_j$
lies in $\ker(H_Z^{(A)})$ (since $G_j$ is a row and the inner product
is $1+1=0$) and is not in $\rowsp(H_X^{(A)} = M^T)$ (minimum Z-stabiliser
weight 4). Exhaustive check confirms 18 weight-2 X-logicals.
\end{proof}

\begin{remark}[Why Dual~A fixes what Code~I cannot]
The rank-gap mechanism (\cref{thm:rankgap}) explains the improvement.
In Code~I, the $G_j$ groups are X-logicals (weight-4, outside
$\rowsp(M^T)$), and the all-ones vector creates 66 weight-2 Z-logicals.
In Dual~A, the $G_j$ groups become Z-\emph{stabilisers}: they are in
$\rowsp(H_Z^{(A)})$ by construction.
This removes the weight-2 and weight-4 Z-logicals, raising $d_Z$ from
2 to 6 at the cost of $d_X$ falling from 4 to 2.
The duality exchanges the roles of $G_j$: logical~$\to$~stabiliser
and stabiliser~$\to$~logical.
\end{remark}

\begin{proposition}[Dual~A code-capacity performance]\label{prop:dualA_perf}
\NEW\; Under independent depolarising noise (20000 trials):
\begin{center}
\small
\begin{tabular}{ccc}
\toprule
$p$ & $P_{\log}(X)$ & $P_{\log}(Z)$ \\
\midrule
0.001 & 0.0120 & 0.0000 \\
0.005 & 0.0480 & 0.0004 \\
0.010 & 0.0807 & 0.0021 \\
0.050 & 0.2767 & 0.0367 \\
\bottomrule
\end{tabular}
\end{center}
Dual~A's Z-threshold exceeds $p=0.20$; its X-threshold is below $p=0.001$.
At $p=0.001$, Dual~A corrects essentially all Z-errors while
simultaneously detecting (but not correcting) X-errors.
\end{proposition}

\subsection{Dual~II: The Z-Dominant Balanced Dual}

\begin{construction}[Dual~II]\label{con:dualII}
\NEW\;
$H_X^{(\mathrm{II}')} := M^T \in \GFtwo^{21\times12}$
(21 plaquettes as X-type checks).
$H_Z^{(\mathrm{II}')} := H_X^{(II)} \in \GFtwo^{4\times12}$
(Code~II's X-check matrix promoted to Z-checks).
\end{construction}

\begin{theorem}[CSS commutativity of Dual~II]\label{thm:dualII_comm}
$H_X^{(\mathrm{II}')} (H_Z^{(\mathrm{II}')})^T = \bm{0}$ over $\GFtwo$.
\end{theorem}
\begin{proof}
Entry $(p, i)$ is $|\text{plaquette } p \cap \text{support of row } i
\text{ of } H_X^{(II)}| \bmod 2$.
Each row of $H_X^{(II)}$ has weight 6 and lies in $\ker(M^T)$, meaning
every plaquette shares an even number of qubits with it.
Therefore all entries are 0 mod 2. 
\end{proof}

\begin{theorem}[Parameters of Dual~II]\label{thm:dualII}
\NEW\; Dual~II has parameters $[[12, 1, (d_X,d_Z)=(3,4)]]$.
\end{theorem}
\begin{proof}
$k = 12 - 7 - 4 = 1$.
$d_Z = 4$: exhaustive verification finds 3 weight-4 Z-logicals and no
lighter ones (the weight-3 Z-logicals of Code~II are now stabilised by
rows of $H_Z^{(\mathrm{II}')} = H_X^{(II)}$).
$d_X = 3$: exhaustive verification finds 16 weight-3 X-logicals.
\end{proof}

\begin{proposition}[Dual~II code-capacity performance]\label{prop:dualII_perf}
\NEW\;
At $p = 0.01$ (20000 trials), Dual~II achieves
$P_{\log}(Z) = 0.0013$: the lookup decoder corrects all weight-1
Z-errors exactly ($d_Z = 4$ with 12 distinct weight-7 single-qubit
Z-error syndromes, exhaustively verified).
For dephasing-dominated hardware where Z-errors vastly outnumber X-errors
(e.g.\ superconducting qubits before dynamical decoupling), Dual~II is
strictly preferable to Code~II.
\end{proposition}

\begin{remark}[The duality is not self-dual]
Code~II and Dual~II have the same parameters up to the exchange
$(d_X,d_Z) = (4,3) \leftrightarrow (3,4)$.
They are not equivalent codes: Code~II has 16 weight-3 Z-logicals and
Dual~II has 16 weight-3 X-logicals.
The duality $\Phi$ maps one to the other but is not an automorphism.
\end{remark}

%------------------------------------------------------------------
\section{Two-Stage Combined Protocol}\label{sec:twostage}
%------------------------------------------------------------------

\subsection{Protocol definition}

\begin{construction}[Two-Stage Protocol]\label{con:twostage}
\NEW\;
Using the same 21 null plaquettes, perform two measurement rounds:

\emph{Stage~1 (Lorentzian/Z-basis):}
Measure each plaquette as a 4-body $Z\otimes Z\otimes Z\otimes Z$
operator. This is equivalent to the Z-check round of Code~I.
Decode the resulting syndrome using the 21-row lookup decoder to
correct X-errors ($d_X=4$, all weight-1 X-errors correctable).

\emph{Stage~2 (Euclidean/X-basis):}
Measure each plaquette as a 4-body $X\otimes X\otimes X\otimes X$
operator. This is equivalent to the X-check round of Dual~A.
Decode the resulting syndrome to correct Z-errors
($d_Z=6$, all weight-1 and weight-2 Z-errors correctable).

The logical space is that of Dual~A with $k_{\mathrm{eff}} = 2$.
\end{construction}

\subsection{Theoretical guarantees}

\begin{theorem}[Two-Stage Correction Distances]\label{thm:twostage}
\NEW\; The two-stage protocol corrects:
\begin{itemize}[noitemsep]
\item All weight-1 X-errors (from Stage~1, $d_X=4$ of the primal code).
\item All weight-1 and weight-2 Z-errors (from Stage~2, $d_Z=6$ of Dual~A).
\end{itemize}
No additional hardware is required beyond the 21 plaquette measurements,
which are measured twice per logical cycle in alternating Pauli bases.
\end{theorem}

\subsection{Numerical results}

\begin{proposition}[Two-Stage Protocol Performance]\label{prop:twostage_perf}
\NEW\;
Under code-capacity independent depolarising noise (20000 trials):
\begin{center}
\small
\begin{tabular}{cccccc}
\toprule
$p$ & Stage~1 $P_{\log}(X)$ & Stage~2 $P_{\log}(Z)$ & Combined & Code~II & Dual~A \\
\midrule
0.001 & 0.0003 & 0.0000 & 0.0000 & 0.0000 & 0.0120 \\
0.005 & 0.0010 & 0.0000 & 0.0019 & 0.0018 & 0.0480 \\
0.010 & 0.0030 & 0.0020 & 0.0060 & 0.0049 & 0.0807 \\
0.020 & 0.0120 & 0.0103 & 0.0233 & 0.0182 & 0.2041 \\
0.050 & 0.0657 & 0.0293 & 0.0969 & 0.0852 & 0.4311 \\
\bottomrule
\end{tabular}
\end{center}
At $p=0.01$, the combined protocol achieves $P_{\log}=0.006$ with
$k_{\mathrm{eff}}=2$ logical qubits, comparable to Code~II ($P_{\log}=0.005$,
$k=1$) while encoding twice as many logical qubits.
The combined threshold exceeds $p=0.20$.
\end{proposition}

\begin{remark}[Physical cost of the two-stage protocol]
The two-stage protocol measures 42 stabiliser circuits per logical cycle
(21 in Z-basis, 21 in X-basis) rather than the 22 (1 X + 21 Z) of Code~I.
This is a factor-of-2 overhead in syndrome extraction, which is offset
by the $k_{\mathrm{eff}}=2$ logical qubit capacity (versus $k=4$ in Code~I
but with failed Z-correction, or $k=1$ in Code~II).
For neutral atom platforms supporting parallel multi-qubit measurements
in switchable bases, the two-stage protocol is particularly natural.
\end{remark}

%------------------------------------------------------------------
\section{No-Go Theorems}\label{sec:nogo}
%------------------------------------------------------------------

\subsection{CSS Pigeonhole No-Go Theorem}

\begin{lemma}\label{lem:allpairs}
Every weight-2 even vector $e_{q_1} + e_{q_2}$ lies in
$\ker(H_X^{(I)}) \setminus \rowsp(\HZ)$.
\end{lemma}
\begin{proof}
Even weight implies membership in $\ker(H_X^{(I)})$.
The minimum weight in $\rowsp(\HZ)$ is 4; exhaustive check over all
66 pairs confirms none is a stabiliser.
\end{proof}

\begin{theorem}[Pigeonhole No-Go]\label{thm:nogo}
\NOGO\;
$d_Z \geq 3$ with $k \geq 2$ is impossible in any CSS code on this
geometry.
More precisely: starting from any single-row $H_X$ (giving $k=4$ with
the all-ones choice), adding $m$ independent rows to reach $d_Z \geq 3$
forces $k \leq 1$.
\end{theorem}
\begin{proof}
All $\binom{12}{2}=66$ qubit pairs are Z-logicals under any even-weight
single-row $H_X$ (\cref{lem:allpairs}).
An X-row $r$ kills pair $(q_1,q_2)$ iff $r[q_1]\oplus r[q_2]=1$:
a binary partition of 12 qubits.
With $m$ independent rows, at most $2^m$ partition classes.
By the Pigeonhole Principle, $m \leq 3$ always leaves an unkilled pair.
For $m=4$: $k = 12-(1+4)-7=0$.

Tight bounds by exhaustive search:
$m=1$: 30 residual pairs; $m=2$: 12; $m=3$: 4; $m=4$: 0 but $k=0$.
Code~II ($m=3$, $k=1$, $d_Z=3$) is the unique tight solution.
\end{proof}

\begin{remark}
The no-go theorem applies to the \emph{primal} family only.
The dual family evades it by using the plaquettes as X-checks
rather than Z-checks: Dual~A achieves $k=2$ with $d_Z=6$ because
the Pigeonhole argument on Z-error pairs does not apply when
$H_X = M^T$ (the 21-row plaquette matrix).
The constraint that $m \leq 3$ independent rows kill all pairs
is a statement about the structure of $\ker(\HZ)$, not about the
geometry per se.
\end{remark}

\subsection{X-Decoration Equivalence}

\begin{theorem}[X-Decoration Equivalence]\label{thm:nodeco}
\NOGO\;
No non-CSS stabiliser code built by augmenting or replacing rows of $\HZ$
with weight-$\leq6$ non-CSS generators achieves $d_Z \geq 3$.
\end{theorem}
\begin{proof}
For a non-CSS generator $(a|b)$ to commute with all 21 pure-Z plaquettes,
its X-part $a$ must lie in $\ker(\HZ)$.
Exhaustive enumeration shows every weight-$\leq6$ vector in $\ker(\HZ)$
is a valid CSS X-row (the 50 orbit representatives under $B_4$ all have
X-part equal to one of the 27 CSS vectors).
The Z-decoration $b$ (yielding Y Paulis) does not change the binary
partition induced by $a$: pair $(q_1,q_2)$ is killed iff
$a[q_1]\oplus a[q_2]=1$. The Pigeonhole bound applies identically.
\end{proof}

%------------------------------------------------------------------
\section{Structural Connection to the Boundary Sum}\label{sec:lorentz}
%------------------------------------------------------------------

\begin{remark}[Structural parallel (primal)]\label{rem:lorentz_primal}
The companion paper establishes
$n_{\mathrm{eff}} = (12, 0, 0, 0)$ (Theorem~4.2 of~\cite{SchmittCD}).
The temporal component $n_{\mathrm{eff}}^0 = 12$ corresponds to the
weight of the all-ones X-check in Code~I.
The spatial components $n_{\mathrm{eff}}^j = 0$ correspond to the
pairwise cancellation within each $G_j$, which places $G_j$ in
$\ker(\HZ)$ and makes them X-logicals in Code~I.
This is a structural parallel, not a logical derivation:
the boundary sum motivates the all-ones choice but does not force it.
\end{remark}

\begin{remark}[Structural parallel (dual)]\label{rem:lorentz_dual}
Under Euclidean orientation, the boundary sum vanishes:
$n_{\mathrm{eff},\mathrm{Eucl}} = (0, 0, 0, 0)$.
This corresponds precisely to the dual family: without a temporal charge,
there is no all-ones check, and the $G_j$ groups transition from
X-logicals to Z-stabilisers, elevating $d_Z$ from 2 to 6.
\end{remark}

%------------------------------------------------------------------
\section{Symmetry Group}\label{sec:symmetry}
%------------------------------------------------------------------

\begin{proposition}\label{prop:symmetry}
The subgroup of the hyperoctahedral group $B_4$ (order 384) acting
faithfully on the 12 null links has order 96.
It preserves $\rowsp(\HZ)$ for all four codes and permutes the three
weight-4 X-logicals of Code~II (equivalently, the three weight-4
Z-logicals of Dual~II) among themselves.
Under the duality $\Phi$, the order-96 group action is preserved:
the same symmetry acts on both the primal and dual families.
\end{proposition}

%------------------------------------------------------------------
\section{Hardware Considerations}\label{sec:hardware}
%------------------------------------------------------------------

\begin{proposition}
The qubit interaction graph is $K_{12}$ (66 edges, non-planar) for
all four codes: they share the same physical qubit layout.
All-to-all connectivity is required.
\end{proposition}

\paragraph{Target platforms.}
Neutral atom arrays (QuEra, Pasqal) and trapped-ion QCCD (IonQ,
Quantinuum) provide the required connectivity.
Neutral atoms natively convert leakage to detectable erasures~\cite{Wu2022};
Code~I's $d_Z=2$ then allows post-selection.
The two-stage protocol is particularly natural on neutral atom platforms,
which support arbitrary multi-qubit observables in switchable Pauli bases
without additional hardware cost.
Code~II and Dual~II each require 37 total qubits (12 data + 25 ancilla),
fitting within current device capacities~\cite{Evered2023}.
Dual~A requires 36 total qubits (12 data + 3 Z-ancilla + 21 X-ancilla).

\paragraph{Asymmetric error models.}
For superconducting qubits before dynamical decoupling, $T_2 \ll T_1$
implies Z-errors dominate.
Dual~II $[[12,1,(3,4)]]$ is strictly preferable to Code~II $[[12,1,(4,3)]]$
in this regime: $d_Z=4$ versus $d_Z=3$ at equal $n$, $k$.

%------------------------------------------------------------------
\section{Circuit-Level Threshold (Code~II)}\label{sec:circuit}
%------------------------------------------------------------------

We simulate one syndrome round of Code~II under the standard
circuit-level depolarising model: each two-qubit (CNOT) gate applies
a random non-identity Pauli with probability $p$; state preparation
and measurement each have error probability $p$.
Syndromes are extracted over $T=3$ rounds and decoded with majority
vote followed by single-error lookup.

\begin{proposition}[Circuit depths]\label{prop:depths}
Code~II requires X-circuit depth 11 and Z-circuit depth 16.
Code~I requires X-circuit depth 4 (three disjoint parallel groups)
and Z-circuit depth 16.
\end{proposition}

\begin{proposition}[Circuit-level logical error rates]\label{prop:circuit}
Under the above model with $T=3$ syndrome rounds (10000 trials per point):
\begin{center}
\small
\begin{tabular}{ccc}
\toprule
$p$ & Code~II $P_{\log}$ & Note \\
\midrule
0.001 & 0.0088 & current best hardware \\
0.003 & 0.0189 & \\
0.005 & 0.0290 & \\
0.008 & 0.0589 & \\
0.010 & 0.0835 & NISQ target \\
0.015 & 0.1556 & \\
0.020 & 0.2335 & \\
0.030 & 0.4136 & \\
0.050 & 0.7335 & \\
\bottomrule
\end{tabular}
\end{center}
Code~II has circuit-level threshold $p_c \approx 3.5\%$, exceeding the
$1\%$ NISQ hardware target.
Code~I has substantially elevated $P_{\log}$ at all tested $p$
(approximately $2.8\text{--}3.5\times$ that of Code~II), because the single
all-ones X-check maps all 12 single-qubit Z-error syndromes to the same
1-bit syndrome value — the lookup decoder cannot locate the error, and any
wrong correction leaves a weight-2 residual that is a Z-logical.
Code~I is a Z-error-\emph{detecting} code; it requires post-selection or
erasure conversion for autonomous Z-correction.
Circuit-level thresholds for Dual~A and Dual~II are left as open
problems (\cref{sec:open}).
\end{proposition}

%------------------------------------------------------------------
\section{Discussion}\label{sec:discussion}
%------------------------------------------------------------------

\subsection{Summary of results}

\begin{enumerate}[label=(\roman*),noitemsep]
\item \textbf{Lorentzian CSS Duality} (\cref{thm:duality}): a geometric
  duality generated by orientation reversal maps
  Code~I~$\to$~Dual~A and Code~II~$\to$~Dual~II.
\item \textbf{GF(2) Rank Gap} (\cref{thm:rankgap}): the rank gap
  $\rank(M)_\R - \rank(M)_{\GFtwo} = 1$ is the algebraic cause of
  $d_Z = 2$ in Code~I, provably forced by the Lorentzian temporal
  boundary charge $n_{\mathrm{eff}}^0 = 12$.
\item \textbf{Dual~A} $[[12,2,(2,6)]]$: corrects all weight-1 and
  weight-2 Z-errors; CSS commutativity proven from the $2\FS+2\PS$
  plaquette structure.
\item \textbf{Dual~II} $[[12,1,(3,4)]]$: a new code with $d_Z = 4$,
  strictly better than Code~II for dephasing-dominated hardware;
  $P_{\log}(Z) = 0.001$ at $p = 0.01$.
\item \textbf{Isometric syndrome structure} (\cref{prop:isometric}):
  the geometry provides equally rich syndrome information for both
  X- and Z-errors; the two code families each exploit one half.
\item \textbf{Two-stage protocol} (\cref{thm:twostage}): using the same
  21 plaquettes in alternating Pauli bases achieves $P_{\log}=0.006$
  with $k_{\mathrm{eff}}=2$ at $p=0.01$.
\item \textbf{Pigeonhole No-Go} (\cref{thm:nogo}): $d_Z\geq3$ with
  $k\geq2$ is impossible in the primal family; Code~II is tight.
\item \textbf{X-Decoration Equivalence} (\cref{thm:nodeco}):
  weight-$\leq6$ non-CSS codes do not escape the No-Go.
\item \textbf{Role transition}: the all-ones vector transitions from
  X-stabiliser (Code~I) to X-logical (Code~II) to Z-logical (Dual~A).
  The $G_j$ groups transition from X-logicals (Code~I) to
  X-stabilisers (Code~II context) to Z-stabilisers (Dual~A).
\end{enumerate}

\subsection{The code family is complete in the CSS weight-$\leq6$ sector}

By the Pigeonhole No-Go (\cref{thm:nogo}), no $[[12,k\geq2,d_Z\geq3]]$
CSS code exists on this geometry.
The four codes identified here are therefore the \emph{complete} CSS
family within weight~$\leq6$ stabilisers, partitioned by the duality
map $\Phi$ into two conjugate pairs.

\subsection{Open problems}\label{sec:open}

\begin{description}[style=nextline]
\item[(O1) Non-CSS balanced code with $k \geq 2$.]
  Find a $[[12, k{\geq}2, 3]]$ non-CSS stabiliser code using
  weight-$\geq8$ generators guided by the $B_4$ orbit structure.
  This is the only remaining gap above the CSS No-Go bound.
\item[(O2) Circuit-level thresholds for Dual~A and Dual~II.]
  The greedy CNOT schedule for the 21 X-type plaquette checks in
  Dual~A has depth 16 (identical to the Z-circuit of Code~I).
  Simulating the circuit-level threshold under the standard model is
  the natural next verification step.
\item[(O3) Circuit-level threshold for the two-stage protocol.]
  The combined protocol doubles the syndrome circuit depth.
  Characterising the circuit-level $P_{\log}$ and threshold for
  the two-stage measurement sequence is an open problem.
\item[(O4) Transversal gate set.]
  Map the order-96 symmetry action to explicit logical gate
  implementations for all four codes, exploiting the duality structure.
\item[(O5) Asymmetric noise experiments on Dual~II.]
  Benchmark Dual~II against Code~II on superconducting hardware
  with realistic biased noise ($Z$-errors $10\times$ more likely than
  $X$-errors), where Dual~II's $d_Z=4$ advantage should be maximal.
\end{description}

%------------------------------------------------------------------
\section*{Acknowledgements}
%------------------------------------------------------------------

The author thanks three anonymous reviewers of~\cite{SchmittCD} whose
detailed evaluations strengthened the geometric foundations of this work,
and a subsequent reviewer who identified the CW-complex obstruction that
led to the corrected flat-connection analysis.

%------------------------------------------------------------------
\begin{thebibliography}{99}

\bibitem{Fowler2012}
A.~G.~Fowler, M.~Martinis, A.~Whitfield and J.~M.~Martinis,
Phys.\ Rev.\ A \textbf{86} (2012) 032324.

\bibitem{Panteleev2022}
P.~Panteleev and G.~Kalachev,
IEEE Trans.\ Inf.\ Theory \textbf{68} (2022) 7820.

\bibitem{Leverrier2022}
A.~Leverrier and G.~Z\'emor,
Proc.\ FOCS (2022) 872.

\bibitem{Schmitt2026a}
Schmitt, Y.\ (2026).
\emph{Exact Discretisation and Boundary Observables in Lorentzian Causal
Diamonds}. Zenodo. https://doi.org/10.5281/zenodo.19338306
Preprint.

\bibitem{Wu2022}
Y.~Wu, S.~Kolkowitz, S.~Puri and J.~D.~Thompson,
Nat.\ Commun.\ \textbf{13} (2022) 4657.

\bibitem{Sahay2023}
K.~Sahay and B.~Neyenhuis,
PRX Quantum \textbf{4} (2023) 020342.

\bibitem{Evered2023}
S.~J.~Evered et al.,
Nature \textbf{622} (2023) 268.

\bibitem{Ryan-Anderson2021}
C.~Ryan-Anderson et al.,
PRX Quantum \textbf{2} (2021) 030307.

\end{thebibliography}

\end{document}